% ─────────────────────────────────────────────────────────────────────────────
% Abstract
% Paper: "Silent Decisions Are Type Errors: Enforcing AI Accountability
%         via Linear Resource Types"
% Venue: IEEE Computer — Special Issue on AI Governance
% Deadline: 13 April 2026
% Authors: Daniel Lau Pereira Soares
% ─────────────────────────────────────────────────────────────────────────────

\begin{abstract}
As autonomous AI agents increasingly make decisions with legal
consequences---credit, employment, healthcare---regulatory frameworks
such as Brazil's LGPD (Art.~18) and the EU AI Act (Art.~86) demand
strict, auditable accountability.
Current Policy-as-Code (PaC) tools treat accountability as an
asynchronous runtime property, generating post-hoc audit logs.
We show that this architectural decoupling creates a structural
vulnerability: \emph{silent decisions}---acts with legal consequences
that lack a computationally bound, non-repudiable evidence chain.

We introduce \emph{Bound-To-Verdict} (BTV), a framework that shifts
AI accountability from a runtime audit obligation to a compile-time
type guarantee.
Drawing on Girard's Linear Logic~\cite{girard1987linear}, BTV encodes
regulatory obligations as linear resource types.
Under the BTV invariant, a decision \texttt{Verdict} ($V$) cannot
materialize without atomically consuming its forensic
\texttt{EvidenceToken} ($E$) and \texttt{ComplianceToken} ($C$),
formalised as:
\[
  V \multimap (E \otimes C).
\]
We prove the \emph{Constitutional Enclosure Theorem}: in any Safe Rust
program implementing this invariant, the set of materialized silent
decisions $\mathcal{S}$ is empty by construction---an attempt to issue
a decision without consuming its evidence is rejected by the compiler
as a type error, not flagged at runtime.

We provide a reference implementation in Safe Rust
(\texttt{silent-decisions-proof}) to validate empirical viability.
Criterion benchmarks show that this static guarantee operates as a
zero-cost abstraction: the governance layer introduces 152 bytes of
memory overhead and $1.67\,\mu$s of latency for 4\,KB payloads, with
all overhead attributable exclusively to cryptographic primitives
(BLAKE3 + HMAC-SHA256) that any accountable system must perform
regardless of governance architecture.
BTV demonstrates that cryptographic, non-repudiable AI accountability
can be enforced structurally without compromising high-throughput
performance.
\end{abstract}
